% Bảng 1. Nguồn số: facts.md mục "Bổ sung sau phản biện 3" (nguyên nhân lọt v1: Test-I 13 = SL 6, AL 3, PT 4 (3 sửa giá
% trị + 1 đích có trong u); Test-II 27 = SL 14, AL 2, PT 7, US 3, CR 1) và mục Test-4 ("40 ca vượt v1": M3 20, M5 11,
% M4 5, M6 3, M2 1). Tổng theo nguyên nhân: nguồn công khai 25, đích có trong u 5, sửa tham số 10.
\begin{table}[!htbp]
\centering
\caption{Sáu cơ chế tấn công. Cột thứ ba ghi số ca lọt qua \base{} trên hai tập niêm phong trước (\tOne{} và \tTwo,
ba mô hình nền, tổng 40 ca) và nguyên nhân. Cột cuối là thành phần của \ours{} nhắm vào nhóm đó.}\label{tab:mechanisms}
\footnotesize
\setlength{\tabcolsep}{3pt}
\begin{tabularx}{\textwidth}{@{}>{\raggedright\arraybackslash}p{2.45cm}>{\raggedright\arraybackslash}X%
>{\raggedright\arraybackslash}p{4.0cm}l@{}}
\toprule
Cơ chế & Con đường gây hại & Lọt qua \base{} (nguyên nhân) & \ours \\
\midrule
\mER{} chuyển hướng tường minh & Đoạn chèn yêu cầu gửi dữ liệu tới một đích lạ & 0 & \cOne \\
\mCR{} chuyển hướng che giấu & Như trên, nhưng diễn đạt gián tiếp hoặc tách nhỏ đích & 1 (trường đích là một từ thường
  có trong $u$) & \cOne \\
\mSL{} rò rỉ qua sink & Giá trị mật đi vào nội dung gửi tới một đích hợp lệ & 20 (giá trị nằm trong nguồn gắn nhãn
  công khai) & \cTwo \\
\mAL{} rò rỉ qua câu trả lời & Giá trị mật xuất hiện trong câu trả lời cuối & 5 (như trên) & \cTwo \\
\mPT{} sửa tham số & Đổi một giá trị người dùng đã nêu, giữ đích hợp lệ & 11 (10 ca sửa giá trị, 1 ca đích chỉ được
  nhắc tới) & \cThree, \cOne \\
\mUS{} tác dụng phụ ngoài yêu cầu & Gọi sink tới một đích chỉ được nhắc tới hoặc bị cấm & 3 (đích có trong $u$ nhưng
  không được cấp quyền) & \cOne \\
\bottomrule
\end{tabularx}
\end{table}
